% !PRINTED 108
\weberpage{108}
Let us begin, then, with the cubic equations, and understand
by $x_0, x_1, x_2$ the roots of any \gesperrt{Abelian}
equation of the third degree. Since 3 is a prime, the
equation must be cyclic, i.\,e.\ \gesperrt{the cyclic functions
of the quantities} $x_0, x_1, x_2$ \gesperrt{are rational numbers.}

If we denote by $\varrho$ an imaginary third root of unity,
we have the two resolvents
\weq{1}{\begin{aligned}
(\varrho, x_0) &= x_0 + \varrho\, x_1 + \varrho^2 x_2\\
(\varrho^2, x_0) &= x_0 + \varrho^2 x_1 + \varrho\, x_2,
\end{aligned}}
whose cubes, then. must be numbers of the field $R(\varrho)$, and whose product % [sic: Weber prints a full stop after ``also'', although the sentence continues.]
must be a rational number.

Accordingly, denoting by $a, b, c$ rational (integer
or fractional) numbers, we set
\weq{2}{\begin{aligned}
(\varrho, x_0)^3 &= a + b\varrho\\
(\varrho^2, x_0)^3 &= a + b\varrho^2\\
(\varrho, x_0)(\varrho^2, x_0) &= c,
\end{aligned}}
and hence also
\weq{3}{(a+b\varrho)(a+b\varrho^2) = c^3.}

Now in the field $R(\varrho)$ we have, besides the six units
\[ \pm 1,\ \pm\varrho,\ \pm\varrho^2, \]
which can all be represented as powers of $-\varrho$, the
primes $\sqrt{-3} = \varrho - \varrho^2$, the real primes $q$ of the
form $3N+2$, and the two complex factors of the real
primes $p$ of the form $3N+1$. The decomposition of these
primes into the two complex factors $\pi, \pi'$:
\weq{4}{p = \pi\pi'}
we have presented in \S.\,172 of the first volume in the form
\weq{5}{p = \psi_1(\varrho)\,\psi_1(\varrho^2).}
We can therefore set
\weq{6}{\pi = \psi_1(\varrho), \quad \pi' = \psi_1(\varrho^2),}
and have thereby selected, among the various numbers associated
with $\pi$, one determinate number; and these numbers
$\pi, \pi'$ stand in a determinate relation to the cyclotomic
periods of the $p$th roots of unity $\eta$:
\weq{7}{(\varrho, \eta)^3 = p\pi, \quad (\varrho^2, \eta)^3 = p\pi'}
[Bd.\ I, \S.\,172, (5)].
